En un centre d'estètica disposen d'una màquina de bronzejar amb radiació ultraviolada. Cal canviar-ne un dels tubs fluorescents perquè s'ha fet malbé per l'ús. El tub que cal substituir indica: «Llum UVA 300~nm 20~W 600~mm». Atès que el fabricant de la màquina ha fet fallida, es busca un tub fluorescent compatible. Després de fer-ne una selecció, es tria un llum que pot fer servei. A les especificacions del producte triat hi diu: «Llum UVA 350~nm 20~W 600~mm T8 làmpada fluorescent». Com que els dos tubs consumeixen la mateixa potència, emeten el mateix nombre de fotons.

L'aparell disposa d'un dispositiu de seguretat basat en l'efecte fotoelèctric que apaga el fluorescent quan el nombre d'electrons emesos per unitat de temps és superior a $2,50\times 10^{15}\ \mathrm{s^{-1}}$. Aquest dispositiu està format per una placa de sodi (la funció de treball és 2,40~eV) i, amb els tubs originals, el nombre d'electrons que abandona la superfície de sodi per unitat de temps és $2,00\times 10^{15}\ \mathrm{s^{-1}}$.

\begin{apartat}{1,25}
Determineu l'energia cinètica màxima dels electrons emesos i la intensitat de corrent que abandona la superfície de sodi amb els tubs originals.
\end{apartat}

\begin{apartat}{1,25}
Determineu com afecta el nou tub al funcionament del dispositiu de seguretat.
\end{apartat}

\dades{%
  $1\ \mathrm{eV} = 1,602\times 10^{-19}\ \mathrm{J}$.\\
  $|e| = 1,602\times 10^{-19}\ \mathrm{C}$.\\
  $c = 3,00\times 10^{8}\ \mathrm{m\,s^{-1}}$.\\
  $h = 6,63\times 10^{-34}\ \mathrm{J\,s}$.}
